\documentclass[12pt]{article}

\usepackage[T1]{fontenc}
\usepackage[utf8]{inputenc}
\usepackage[brazil]{babel}
\usepackage{lmodern}
\usepackage{microtype}
\usepackage[a4paper,margin=2.5cm]{geometry}
\usepackage{setspace}
\usepackage{indentfirst}
\usepackage{amsmath}

\setstretch{1.15}

\title{Palestra do Nobel: Inflação e Desemprego \thanks{Texto Traduzido livremente por Luiz Mario Andrade com o auxílio da ferramenta Chat GPT 5.4}}
\author{Milton Friedman}
\date{Junho de 1977}

\begin{document}

\maketitle

\begin{abstract}
Nas últimas décadas, as visões profissionais sobre a relação entre inflação e desemprego passaram por dois estágios e estão agora entrando em um terceiro. O primeiro foi a aceitação de um \textit{trade-off} estável (uma curva de Phillips estável). O segundo foi a introdução das expectativas inflacionárias como uma variável que desloca a curva de Phillips de curto prazo e da taxa natural de desemprego como determinante da localização de uma curva de Phillips vertical de longo prazo. O terceiro é ocasionado pelo fenômeno empírico de uma aparente relação positiva entre inflação e desemprego. O artigo explora a possibilidade de que essa relação possa ser mais do que coincidente.
\end{abstract}

\vspace{1cm}

Quando o Banco da Suécia estabeleceu o prêmio em ciências econômicas em memória de Alfred Nobel em 1968, havia sem dúvida — como ainda há — um ceticismo generalizado entre cientistas e o público em geral sobre a adequação de tratar a economia como paralela à física, química e medicina. Estas são vistas como "ciências exatas" nas quais o conhecimento objetivo, cumulativo e definitivo é possível. A economia e suas ciências sociais correlatas são vistas mais como ramos da filosofia do que como ciência propriamente definida, imersas em valores desde o início porque lidam com o comportamento humano. Não seriam as ciências sociais, nas quais os estudiosos estão analisando o comportamento de si mesmos e de seus semelhantes, que por sua vez observam e reagem ao que os estudiosos dizem, fundamentalmente dependentes de métodos de investigação diferentes das ciências físicas e biológicas? Não deveriam elas ser julgadas por critérios diferentes?

\section{Ciências Sociais e Naturais}

Eu nunca aceitei essa visão. Acredito que ela reflete uma incompreensão tanto do caráter quanto das possibilidades da ciência social, bem como do caráter e das possibilidades da ciência natural. Em ambas, não existe conhecimento substantivo "certo"; apenas hipóteses tentativas que nunca podem ser "provadas", mas que podem apenas falhar em ser rejeitadas, hipóteses nas quais podemos ter mais ou menos confiança, dependendo de características como a amplitude da experiência que abrangem em relação à sua própria complexidade e em relação a hipóteses alternativas, e o número de ocasiões em que escaparam da rejeição possível. Tanto nas ciências sociais quanto nas naturais, o corpo de conhecimento positivo cresce pela falha de uma hipótese tentativa em prever fenômenos que a hipótese professa explicar; pelo remendo dessa hipótese até que alguém sugira uma nova hipótese que incorpore de forma mais elegante ou simples os fenômenos problemáticos, e assim por diante \textit{ad infinitum}. Em ambas, o experimento é às vezes possível, às vezes não (testemunhe a meteorologia). Em ambas, nenhum experimento é jamais completamente controlado, e a experiência frequentemente oferece evidências que são o equivalente a um experimento controlado. Em ambas, não há como ter um sistema fechado autossuficiente ou evitar a interação entre o observador e o observado. O teorema de Gödel na matemática, o princípio da incerteza de Heisenberg na física, a profecia autorrealizável ou autoderrotável nas ciências sociais exemplificam essas limitações.

É claro que as diferentes ciências lidam com diferentes assuntos, têm diferentes corpos de evidência para extrair (por exemplo, a introspecção é uma fonte de evidência mais importante para o social do que para as ciências naturais), encontram diferentes técnicas de análise mais úteis e alcançaram sucesso diferencial na previsão dos fenômenos que estão estudando. Mas tais diferenças são grandes entre, digamos, física, biologia, medicina e meteorologia, tanto quanto entre qualquer uma delas e a economia.

Até mesmo o difícil problema de separar julgamentos de valor de julgamentos científicos não é exclusivo das ciências sociais. Lembro-me bem de um jantar em um colégio da Universidade de Cambridge, onde eu estava sentado entre um colega economista e R. A. Fisher, o grande estatístico matemático e geneticista. Meu colega economista me falou sobre um aluno que ele estava orientando em economia do trabalho, que, em conexão com uma análise do efeito dos sindicatos, comentou: "Bem, certamente o Sr. X (outro economista de uma persuasão política diferente) não concordaria com isso". Meu colega considerou essa experiência como uma acusação terrível contra a economia porque ilustrava a impossibilidade de uma ciência econômica positiva livre de valores. Voltei-me para Sir Ronald e perguntei se tal experiência era de fato exclusiva da ciência social. Sua resposta foi um "não" apaixonado, e ele passou a contar uma história após outra sobre a precisão com que se podia inferir visões em genética a partir de visões políticas.

Um dos meus grandes professores, Wesley C. Mitchell, gravou em mim a razão básica pela qual os estudiosos têm todo incentivo para buscar uma ciência livre de valores, por mais que seus valores e quão fortemente eles possam desejar espalhá-los e promovê-los. Para recomendar um curso de ação para atingir um objetivo, devemos primeiro saber se esse curso de ação de fato promoverá o objetivo. O conhecimento científico positivo que nos permite prever as consequências de um possível curso de ação é claramente um pré-requisito para o julgamento normativo sobre se esse curso de ação é desejável. O Caminho para o Inferno é pavimentado com boas intenções, precisamente por causa da negligência deste ponto bastante óbvio.

Este ponto é particularmente importante na economia. Muitos países ao redor do mundo estão hoje experimentando uma inflação socialmente destrutiva, um desemprego anormalmente alto, mau uso de recursos econômicos e, em alguns casos, a supressão da liberdade humana não porque homens maus buscaram deliberadamente alcançar esses resultados, nem por causa de diferenças de valores entre seus cidadãos, mas por causa de julgamentos errôneos sobre as consequências das medidas governamentais: erros que, pelo menos em princípio, são passíveis de serem corrigidos pelo progresso da ciência econômica positiva.

Em vez de prosseguir com essas ideias no abstrato (discuti as questões metodológicas de forma mais completa em Friedman [1953]), ilustrarei o caráter científico positivo da economia discutindo uma questão econômica específica que tem sido uma grande preocupação da profissão econômica ao longo do período pós-guerra, a saber, a relação entre inflação e desemprego. Esta questão é uma ilustração admirável porque tem sido uma questão política controversa ao longo do período, e ainda assim a mudança drástica que ocorreu nas visões profissionais aceitas foi produzida principalmente pela resposta científica à experiência que contradizia uma hipótese aceita tentativamente — precisamente o processo clássico para a revisão de uma hipótese científica.

Não posso dar aqui uma pesquisa exaustiva do trabalho que foi feito sobre esta questão ou da evidência que levou à revisão da hipótese. Poderei apenas passar pela superfície na esperança de transmitir o sabor desse trabalho e dessa evidência e de indicar os principais itens que requerem investigação adicional.

A controvérsia profissional sobre a relação entre inflação e desemprego tem se entrelaçado com a controvérsia sobre o papel relativo dos fatores monetários, fiscais e outros na influência da demanda agregada. Uma questão lida com a forma como uma mudança na demanda nominal agregada, no entanto produzida, se resolve através de mudanças no emprego e nos níveis de preços; a outra, com os fatores que explicam a mudança na demanda nominal agregada.

As duas questões estão intimamente relacionadas. Os efeitos de uma mudança na demanda nominal agregada sobre o emprego e os níveis de preços podem não ser independentes da fonte da mudança e, inversamente, o efeito da política monetária, fiscal ou outras forças sobre a demanda nominal agregada pode depender de como o emprego e os níveis de preços reagem. Uma análise completa claramente teria que tratar as duas questões conjuntamente. No entanto, há uma medida considerável de independência entre elas. Para uma primeira aproximação, os efeitos sobre o emprego e os níveis de preços podem depender apenas da magnitude da mudança na demanda nominal agregada, não da sua fonte. Em ambas as questões, a opinião profissional hoje é muito diferente do que era logo após a Segunda Guerra Mundial, porque a experiência contradisse hipóteses aceitas tentativamente. Qualquer uma das questões poderia, portanto, servir para ilustrar minha tese principal. Escolhi lidar com apenas uma para manter esta palestra dentro de limites razoáveis. Escolhi fazer dessa questão a relação entre inflação e desemprego, porque a experiência recente me deixa menos satisfeito com a adequação do meu trabalho anterior sobre essa questão do que com a adequação do meu trabalho anterior sobre as forças que produzem mudanças na demanda agregada.

\section{Estágio 1: Curva de Phillips Negativamente Inclinada}

A análise profissional da relação entre inflação e desemprego passou por dois estágios desde o fim da Segunda Guerra Mundial e está agora entrando em um terceiro. O primeiro estágio foi a aceitação de uma hipótese associada ao nome de A. W. Phillips (1958) de que existe uma relação negativa estável entre o nível de desemprego e a taxa de variação de salários — níveis elevados de desemprego sendo acompanhados por salários em queda, níveis baixos de desemprego por salários em ascensão. A variação de salários, por sua vez, estava ligada à variação de preços, permitindo o aumento secular na produtividade e tratando o excesso de preço sobre o custo salarial como dado por um fator de margem (\textit{markup}) aproximadamente constante.

A Figura 1 ilustra esta hipótese, onde segui a prática padrão de relacionar o desemprego diretamente à variação de preços, abreviando a etapa intermediária através dos salários.

\begin{center}
[Figura 1: Curva de Phillips Simples. O eixo vertical é a Taxa de variação de preços (1/P dP/dt). O eixo horizontal é o Desemprego. A curva é descendente da esquerda para a direita, cruzando o eixo horizontal em $U_O$. O ponto A no eixo vertical corresponde a $U_L$ no eixo horizontal.]
\end{center}

Esta relação foi amplamente interpretada como uma relação causal que oferecia um \textit{trade-off} estável aos formuladores de políticas. Eles poderiam escolher uma meta de baixo desemprego, como $U_L$. Nesse caso, teriam que aceitar uma taxa de inflação de $A$. Restaria o problema de escolher as medidas (monetárias, fiscais, talvez outras) que produziriam o nível de demanda nominal agregada necessário para alcançar $U_L$, mas se isso fosse feito, não haveria necessidade de preocupação em manter essa combinação de desemprego e inflação. Alternativamente, os formuladores de políticas poderiam escolher uma taxa de inflação baixa ou até mesmo deflação como seu objetivo. Nesse caso, teriam que se conformar com um desemprego mais alto: $U_O$ para inflação zero, $U_H$ para deflação.

Economistas então se ocuparam em tentar extrair a relação retratada na figura 1 a partir de evidências para diferentes países e períodos, para eliminar o efeito de perturbações estranhas, para esclarecer a relação entre variação salarial e variação de preços, e assim por diante. Além disso, exploraram ganhos e perdas sociais decorrentes da inflação de um lado e do desemprego do outro, a fim de facilitar a escolha do \textit{trade-off} "correto".

Infelizmente para esta hipótese, evidências adicionais falharam em conformar-se a ela. As estimativas empíricas da relação da curva de Phillips não eram satisfatórias. Mais importante, a taxa de inflação que parecia ser consistente com um nível especificado de desemprego não permaneceu fixa: nas circunstâncias do período pós-Segunda Guerra Mundial, quando governos em todos os lugares buscavam promover o "pleno emprego", ela tendia a subir em um país após outro ao longo do tempo e a variar bruscamente entre os países. Olhado de outra forma, taxas de inflação que anteriormente tinham sido associadas a baixos níveis de desemprego eram experimentadas juntamente com altos níveis de desemprego. O fenômeno da inflação alta e do desemprego alto simultâneos forçou-se cada vez mais ao aviso público e profissional, recebendo o rótulo nada amável de "estagflação".

Alguns de nós fomos céticos desde o início sobre a validade de uma curva de Phillips estável, primariamente por fundamentos teóricos e não empíricos (Friedman 1966a, 1966b, 1968a, 1968b). O que importava para o emprego, argumentamos, não eram salários em dólares ou libras ou coroas, mas salários reais — o que os salários comprariam em bens e serviços. Baixo desemprego significaria, de fato, pressão por um salário real mais alto — mas salários reais poderiam ser mais altos mesmo se os salários nominais fossem mais baixos, desde que os preços fossem ainda mais baixos. Da mesma forma, o desemprego elevado significaria pressão por um salário real mais baixo — mas os salários reais poderiam ser mais baixos mesmo se os salários nominais fossem mais altos, desde que os preços subissem ainda mais rápido.

Não há necessidade de assumir uma curva de Phillips estável para explicar a tendência aparente de uma aceleração da inflação para reduzir o desemprego. Isso pode ser explicado pelo impacto de mudanças \textit{antecipadas} na demanda nominal em mercados caracterizados por compromissos de longo prazo (implícitos ou explícitos) com relação a capital e trabalho. Compromissos trabalhistas de longo prazo podem ser explicados pelo custo de aquisição de informações pelos empregadores sobre os empregados e pelos empregados sobre oportunidades alternativas de emprego, além do capital humano específico que torna o valor de um empregado para um empregador específico crescer ao longo do tempo e exceder seu valor para outros empregadores potenciais.

Apenas surpresas importam. Se todos antecipassem que os preços subiriam, digamos, 20\% ao ano, então essa antecipação seria incorporada em contratos futuros de salários (e outros), os salários reais se comportariam precisamente como se todos antecipassem que os preços não subiriam, e não haveria razão para que a taxa de inflação de 20\% fosse associada a um nível de desemprego diferente de uma taxa zero. Uma mudança não antecipada é muito diferente, especialmente na presença de compromissos de longo prazo — eles mesmos em parte resultado do conhecimento imperfeito cujo efeito eles aumentam e espalham ao longo do tempo. Compromissos de longo prazo significam, primeiro, que não há compensação instantânea de mercado (como nos mercados de alimentos perecíveis), mas apenas um ajuste defasado de ambos os preços e quantidades a mudanças na demanda ou oferta (como no mercado de aluguel de casas); segundo, que os compromissos assumidos dependem não apenas dos preços observáveis atuais, mas também dos preços esperados para prevalecer ao longo do prazo do compromisso.

\section{Estágio 2: Hipótese da Taxa Natural}

Prosseguindo nestas linhas, nós (em particular, E. S. Phelps [1967, 1970] e eu [1968b]) desenvolvemos uma hipótese alternativa que distinguia entre os efeitos de curto e longo prazo de mudanças não antecipadas na demanda nominal agregada. Comece a partir de alguma posição inicial estável e deixe que haja, por exemplo, uma aceleração não antecipada da demanda nominal agregada. Isso chegará a cada produtor como um aumento inesperadamente favorável na demanda por seu produto. Em um ambiente no qual mudanças estão sempre ocorrendo na demanda relativa por diferentes bens, ele não saberá se essa mudança é especial para ele ou generalizada. Será racional para ele interpretá-la como sendo pelo menos parcialmente especial e reagir a ela procurando produzir mais para vender ao que ele agora percebe ser um preço de mercado mais alto do que o esperado para a produção futura. Ele estará disposto a pagar salários nominais mais altos do que estava disposto a pagar antes para atrair trabalhadores adicionais. O salário real que importa para ele é o salário em termos do preço de seu produto, e ele percebe esse preço como sendo mais alto do que antes. Um salário nominal mais alto pode, portanto, significar um salário real mais baixo conforme percebido por ele.

Para os trabalhadores, a situação é diferente: o que importa para eles é o poder de compra dos salários, não sobre o bem específico que produzem, mas sobre todos os bens em geral. Tanto eles quanto seus empregadores tendem a ajustar mais lentamente sua percepção dos preços em geral — porque é mais caro adquirir informações sobre isso — do que sua percepção do preço do bem específico que produzem. Como resultado, um aumento nos salários nominais pode ser percebido pelos trabalhadores como um aumento nos salários reais e, portanto, provocar um aumento da oferta ao mesmo tempo em que é percebido pelos empregadores como uma queda nos salários reais e, portanto, provoca um aumento da oferta de empregos. Expressos em termos da média dos preços futuros percebidos, os salários reais são mais baixos; em termos do preço médio futuro percebido, os salários reais são mais altos.

Mas esta situação é temporária: deixe que a taxa mais alta de crescimento da demanda nominal e dos preços continue, e as percepções se ajustarão à realidade. Quando o fizerem, o efeito inicial desaparecerá e será então revertido por um tempo, conforme trabalhadores e empregadores se veem presos em contratos inadequados. No final, o emprego voltará ao nível que prevalecia antes da aceleração não antecipada assumida na demanda nominal agregada.

Esta hipótese alternativa é retratada na figura 2. Cada curva negativamente inclinada é uma curva de Phillips como a da figura 1, exceto que é para uma taxa de inflação antecipada ou percebida específica, definida como a taxa média percebida de variação de preços, \textit{não} a média das taxas percebidas de variação de preços individuais (a ordem das curvas seria invertida para o segundo conceito). Comece no ponto E e deixe a taxa de inflação por qualquer motivo mover-se de A para B e permanecer lá. O desemprego inicialmente declinaria para $U_L$ no ponto F, movendo-se ao longo da curva definida para uma taxa antecipada de inflação $[(1/P)(dP/dt)]^*$ de A. À medida que as antecipações se ajustam, a curva de Phillips de curto prazo se moveria para cima, finalmente para a curva definida para uma taxa antecipada de inflação de B. Concomitantemente, o desemprego se moveria gradualmente de F para G. (Para uma discussão mais completa, veja Friedman [1976], cap. 12.)

\begin{center}
[Figura 2: Curva de Phillips ajustada por expectativas. O eixo vertical é a Taxa de inflação (1/P dP/dt). O eixo horizontal é o Desemprego. Há uma linha vertical no nível $U_N$ de desemprego. Várias curvas descendentes cruzam esta linha. O ponto E está em ($U_N$, A). O ponto F está em ($U_L$, B) em uma curva inferior. O ponto G está em ($U_N$, B) em uma curva superior.]
\end{center}

Esta análise é, obviamente, simplificada. Supõe uma única mudança não antecipada, enquanto, é claro, há um fluxo contínuo de mudanças não antecipadas; não lida explicitamente com defasagens, ou com \textit{overshooting}, ou com o processo de formação de antecipações. Mas destaca os pontos-chave: o que importa não é a inflação \textit{per se}, mas a inflação não antecipada; não há \textit{trade-off} estável entre inflação e desemprego; existe uma "taxa natural de desemprego" ($U_N$) que é consistente com as forças reais e com percepções precisas; o desemprego pode ser mantido abaixo desse nível apenas por uma inflação acelerada; ou acima dele apenas por uma deflação acelerada.

A "taxa natural de desemprego", um termo que introduzi para fazer paralelo à "taxa natural de juros" de Knut Wicksell, não é uma constante numérica, mas depende de fatores "reais" em oposição aos monetários — a eficácia do mercado de trabalho, a extensão da competição ou monopólio, as barreiras ou encorajamentos ao trabalho em várias ocupações, e assim por diante.

Por exemplo, a taxa natural claramente tem subido nos Estados Unidos por duas razões principais. Primeiro, mulheres, adolescentes e trabalhadores de tempo parcial têm constituído uma fração crescente da força de trabalho. Esses grupos são mais móveis no emprego do que outros trabalhadores, entrando e saindo do mercado de trabalho, trocando de emprego com mais frequência. Como resultado, tendem a experimentar taxas médias de desemprego mais altas. Segundo, o seguro-desemprego e outras formas de assistência a pessoas desempregadas tornaram-se mais disponíveis para mais categorias de trabalhadores e tornaram-se mais generosos em duração e valor. Os trabalhadores que perdem seus empregos estão sob menos pressão para procurar outro trabalho e tendem a esperar mais tempo na esperança, geralmente cumprida, de serem chamados de volta ao seu emprego anterior, e podem ser mais seletivos nas alternativas que consideram. Além disso, a disponibilidade do seguro-desemprego torna mais atraente entrar na força de trabalho em primeiro lugar e, portanto, pode ter estimulado tanto o crescimento que ocorreu na força de trabalho como uma porcentagem da população quanto sua composição em mudança.

Os determinantes da taxa natural de desemprego merecem uma análise muito mais completa tanto para os Estados Unidos quanto para outros países. O mesmo ocorre com o significado dos números registrados de desemprego e a relação entre os números registrados e a taxa natural. Essas questões são todas de extrema importância para a política pública. No entanto, são questões secundárias para o meu propósito limitado atual.

A conexão entre o estado do emprego e o nível de eficiência ou produtividade de uma economia é outro tópico que é de importância fundamental para a política pública, mas é uma questão secundária para o meu propósito atual. Existe uma tendência de tomar como certo que um alto nível de desemprego registrado é evidência de uso ineficiente de recursos e, inversamente. Esta visão está seriamente errada. Um baixo nível de desemprego pode ser um sinal de uma economia de recrutamento forçado que está usando seus recursos de forma ineficiente e está induzindo trabalhadores a sacrificar o lazer por bens que eles valorizam menos do que o lazer sob a crença equivocada de que seus salários reais serão mais altos do que provam ser. Ou uma taxa natural baixa de desemprego pode refletir arranjos institucionais que inibem a mudança. Uma economia altamente estática e rígida pode ter um lugar fixo para todos, enquanto uma economia dinâmica e altamente progressiva, que oferece oportunidades em constante mudança e promove a flexibilidade, pode ter uma taxa natural de desemprego elevada. Para ilustrar como a mesma taxa pode corresponder a condições muito diferentes: tanto o Japão quanto o Reino Unido tiveram taxas médias baixas de desemprego de, digamos, 1950 a 1970, mas o Japão experimentou um crescimento rápido, o Reino Unido, estagnação.

A hipótese da "taxa natural" ou "aceleracionista" ou "curva de Phillips ajustada por expectativas" — como tem sido variadamente designada — é agora amplamente aceita pelos economistas, embora de forma alguma universalmente. Alguns ainda se apegam à curva de Phillips original; mais reconhecem a diferença entre as curvas de curto e longo prazo, mas consideram até mesmo a curva de longo prazo como negativamente inclinada, embora de forma mais acentuada do que as curvas de curto prazo; alguns substituem uma relação estável entre a aceleração da inflação e o desemprego por uma relação estável entre inflação e desemprego — cientes de, mas não preocupados com a possibilidade de que a mesma lógica que os levou a uma segunda derivada os levaria a derivadas ainda mais altas.

Muita pesquisa econômica atual é dedicada a explorar vários aspectos deste segundo estágio — a dinâmica do processo, a formação de expectativas e o tipo de política sistemática, se houver, que pode ter um efeito previsível sobre magnitudes reais. Podemos esperar um progresso rápido nessas questões. (Mencão especial deve ser feita ao trabalho sobre "expectativas racionais", especialmente as contribuições seminais de John Muth, Robert Lucas e Thomas Sargent; veja Muth [1961], Gordon [1976].)

\section{Estágio 3: Uma Curva de Phillips Positivamente Inclinada?}

Embora o segundo estágio esteja longe de ter sido totalmente explorado, muito menos totalmente absorvido na literatura econômica, o curso dos eventos já está produzindo uma mudança para um terceiro estágio. Nos últimos anos, a inflação mais alta foi frequentemente acompanhada por um desemprego mais alto, não mais baixo, especialmente por períodos de vários anos de duração. Uma curva de Phillips estatística simples para tais períodos parece ser positivamente inclinada, não vertical. O terceiro estágio é direcionado a acomodar este fenômeno empírico aparente. Para fazê-lo, suspeito que será necessário incluir na análise a interdependência da experiência econômica e dos desenvolvimentos políticos. Terá que tratar pelo menos alguns fenômenos políticos não como variáveis independentes — como variáveis exógenas no jargão econométrico — mas como sendo eles mesmos determinados por eventos econômicos — como variáveis endógenas (Gordon 1975b). O terceiro estágio será, acredito, grandemente influenciado por um terceiro grande desenvolvimento — a aplicação da análise econômica ao comportamento político, um campo no qual o trabalho pioneiro também foi feito por Stigler e Becker, bem como por Kenneth Arrow, Duncan Black, Anthony Downs, James Buchanan, Gordon Tullock e outros.

A aparente relação positiva entre inflação e desemprego tem sido uma fonte de grande preocupação para os formuladores de políticas governamentais. Permitam-me citar uma frase de um discurso recente do Primeiro-Ministro Callaghan da Grã-Bretanha: "Costumávamos pensar que era possível gastar dinheiro para sair de uma recessão e aumentar o emprego cortando impostos e aumentando os gastos do governo. Digo-lhes com toda a franqueza que essa opção não existe mais e que, na medida em que existiu, funcionou apenas injetando doses maiores de inflação na economia, seguidas por níveis mais elevados de desemprego no passo seguinte. Essa é a história dos últimos 20 anos" (discurso na Conferência do Partido Trabalhista, 28 de setembro de 1976).

A mesma visão é expressa em um livro branco do governo canadense: "A inflação contínua, particularmente na América do Norte, tem sido acompanhada por um aumento nas taxas de desemprego medidas" ("The Way Ahead: A Framework for Discussion", Governo do Canadá, Working Paper, outubro de 1976).

Estas são declarações notáveis, pois contradizem diretamente as políticas adotadas por quase todos os governos ocidentais ao longo do período pós-guerra.

\subsection{Algumas Evidências}

Evidências mais sistemáticas para as últimas duas décadas são dadas na tabela 1 e nas figuras 3 e 4, que mostram as taxas de inflação e desemprego em sete países industrializados ao longo das últimas duas décadas. De acordo com as médias de 5 anos na tabela 1, a taxa de inflação e o nível de desemprego moveram-se em direções opostas — o resultado esperado da curva de Phillips simples — em cinco de sete países entre os dois primeiros quinquênios (1956–60, 1961–65); em apenas quatro de sete países entre o segundo e o terceiro quinquênio (1961–65 e 1966–70); e em apenas um de sete países entre os dois quinquênios finais (1966–70 e 1970–75). E até a única exceção — Itália — não é uma exceção real. É verdade que o desemprego médio foi um pouco menor de 1971 a 1975 do que nos 5 anos anteriores, apesar de uma taxa de inflação mais do que triplicada. No entanto, desde 1973, tanto a inflação quanto o desemprego aumentaram acentuadamente.

As médias para todos os sete países plotadas na figura 3 mostram ainda mais claramente a mudança de uma curva de Phillips simples negativamente inclinada para uma positivamente inclinada. As duas curvas movem-se em direções opostas entre os dois primeiros quinquênios; na mesma direção a partir de então.

Os dados anuais na figura 4 contam uma história similar, embora mais confusa. Nos primeiros anos, há uma ampla variação na relação entre preços e desemprego, variando de essencialmente nenhuma relação, como na Itália, a uma relação negativa razoavelmente clara de ano para ano, como no Reino Unido e nos Estados Unidos. Em anos recentes, no entanto, França, Estados Unidos, Reino Unido, Alemanha e Japão mostram um aumento marcante tanto na inflação quanto no desemprego — embora para o Japão o aumento no desemprego seja muito menor em relação ao aumento na inflação do que nos outros países, refletindo o significado diferente de desemprego no ambiente institucional diferente do Japão. Apenas Suécia e Itália falham em se conformar ao padrão geral.

É claro que estes dados são, no máximo, sugestivos. Não temos realmente sete corpos independentes de dados. Influências internacionais comuns afetam todos os países, de modo que multiplicar o número de países não multiplica proporcionalmente a quantidade de evidência. Em particular, a crise do petróleo atingiu todos os sete países ao mesmo tempo. Qualquer que tenha sido o efeito que a crise teve na taxa de inflação, ela perturbou diretamente o processo produtivo e tendeu a aumentar o desemprego. Tais aumentos dificilmente podem ser atribuídos à aceleração da inflação que os acompanhou; no máximo, os dois poderiam ser considerados como, pelo menos em parte, o resultado comum de uma terceira influência (Gordon 1975a).

Tanto os dados quinquenais quanto os anuais mostram que a crise do petróleo não pode explicar totalmente o fenômeno descrito graficamente por Callaghan. Já antes do quadruplicar dos preços do petróleo em 1973, a maioria dos países mostrava uma associação clara de inflação crescente e desemprego crescente. Mas isso também pode refletir forças independentes em vez da influência da inflação no desemprego. Por exemplo, as mesmas forças que têm aumentado a taxa natural de desemprego nos Estados Unidos podem ter operado em outros países e podem ser responsáveis por sua tendência de aumento do desemprego, independentemente das consequências da inflação.

Apesar destas ressalvas, os dados sugerem fortemente que, pelo menos em alguns países, dos quais Grã-Bretanha, Canadá e Itália podem ser os melhores exemplos, a inflação crescente e o desemprego crescente foram mutuamente reforçadores, em vez de os efeitos separados de causas distintas. Os dados não são inconsistentes com a afirmação mais forte de que, em todos os países industrializados, taxas mais elevadas de inflação têm alguns efeitos que, por um tempo, contribuem para um desemprego mais elevado. O restante deste artigo é dedicado a uma exploração preliminar do que alguns desses efeitos podem ser.

\subsection{Uma Hipótese Tentativa}

Conjecturo que uma elaboração modesta da hipótese da taxa natural é tudo o que se requer para explicar uma relação positiva entre inflação e desemprego, embora, é claro, tal relação positiva também possa ocorrer por outras razões. Assim como a hipótese da taxa natural explica uma curva de Phillips negativamente inclinada em períodos curtos como um fenômeno temporário que desaparecerá à medida que os agentes econômicos ajustam suas expectativas à realidade, assim uma curva de Phillips positivamente inclinada em períodos um pouco mais longos pode ocorrer como um fenômeno de transição que desaparecerá à medida que os agentes econômicos ajustam não apenas suas expectativas, mas seus arranjos institucionais e políticos a uma nova realidade. Quando isso for alcançado, acredito que — como a hipótese da taxa natural sugere — a taxa de desemprego será amplamente independente da taxa média de inflação, embora a eficiência da utilização de recursos possa não ser. Inflação alta não precisa significar nem desemprego anormalmente alto, nem anormalmente baixo. No entanto, os arranjos institucionais e políticos que a acompanham, seja como relíquias da história anterior ou como produtos da própria inflação, são propensos a se revelarem antitéticos ao uso mais produtivo dos recursos empregados — um caso especial da distinção entre o estado de emprego e a produtividade de uma economia referida anteriormente.

A experiência em muitos países da América Latina que se ajustaram a taxas de inflação cronicamente altas — experiência que foi analisada perceptivamente por alguns dos meus colegas, particularmente Arnold Harberger (1967) e Larry Sjaastad (1974) — é consistente, acredito, com esta visão.

Na versão da hipótese da taxa natural resumida na figura 2, a curva vertical é para taxas alternativas de inflação totalmente antecipada. Qualquer que seja essa taxa — seja ela negativa, zero ou positiva — ela pode ser incorporada em cada decisão se for totalmente antecipada. A uma inflação antecipada de 20 por cento ao ano, por exemplo, contratos de salários de longo prazo forneceriam um aumento em cada ano que subiria em relação ao salário de inflação zero em exatamente 20 por cento ao ano; empréstimos de longo prazo teriam uma taxa de juros 20 por cento superior à taxa de inflação zero ou um principal que seria aumentado em 20 por cento ao ano; e assim por diante — em suma, o equivalente a uma indexação total de todos os contratos. A alta taxa de inflação teria alguns efeitos reais, ao alterar os saldos de caixa desejados, por exemplo, mas não precisaria alterar a eficiência dos mercados de trabalho, ou a duração ou termos dos contratos de trabalho e, portanto, não precisaria mudar a taxa natural de desemprego.

Esta análise supõe implicitamente, primeiro, que a inflação é estável ou pelo menos não mais variável a uma taxa elevada do que a uma taxa baixa — caso contrário, é improvável que a inflação seja tão plenamente antecipada a taxas elevadas quanto a taxas baixas de inflação; segundo, que a inflação é, ou pode ser, aberta, com todos os preços livres para se ajustarem à taxa mais elevada, de modo que os ajustes de preços relativos sejam os mesmos com uma inflação de 20 por cento e com uma inflação zero; terceiro, realmente uma variante do segundo ponto, que não há obstáculos à indexação de contratos.

Em última análise, se a inflação a uma taxa média de 20 por cento ao ano prevalecesse por muitas décadas, essas exigências poderiam chegar razoavelmente perto de serem atendidas, razão pela qual estou inclinado a reter a curva de Phillips vertical de longo prazo. Mas quando um país inicialmente se move para taxas de inflação mais elevadas, essas exigências serão sistematicamente afastadas. E tal período de transição pode muito bem se estender por décadas.

Considere, em particular, os Estados Unidos e o Reino Unido. Por dois séculos antes da Segunda Guerra Mundial para o Reino Unido, e um século e meio para os Estados Unidos, os preços variaram em torno de um nível aproximadamente constante, mostrando aumentos substanciais em tempos de guerra, e declínios subsequentes para níveis de pré-guerra aproximados. O conceito de um nível de preço "normal" estava profundamente imerso nas instituições financeiras e outras de ambos os países e nos hábitos e atitudes de seus cidadãos.

No período imediato pós-Segunda Guerra Mundial, a experiência anterior era amplamente esperada para recorrer. O fato de que a inflação do pós-guerra foi sobreposta à inflação de guerra; no entanto, a expectativa em ambos os países era de deflação. Levou muito tempo para que o medo da deflação do pós-guerra se dissipasse — se é que já se dissipou — e ainda mais tempo antes que as expectativas começassem a se ajustar à mudança fundamental no sistema monetário. Esse ajuste ainda está longe de estar completo (Klein 1975).

De fato, não sabemos o que um ajuste completo consistirá. Não podemos saber agora se as nações industrializadas retornarão ao padrão pré-Segunda Guerra Mundial de um nível de preço estável a longo prazo, ou se se moverão em direção ao padrão latino-americano de taxas de inflação cronicamente altas — como ocorreu recentemente no Chile e na Argentina (Harberger 1976) — ou se passarão por uma mudança econômica e política ainda mais radical levando a uma resolução ainda diferente da presente situação ambígua.

Esta incerteza — ou mais precisamente, as circunstâncias que produzem esta incerteza — leva a desvios sistemáticos das condições exigidas para uma curva de Phillips vertical.

O desvio mais fundamental é que uma taxa de inflação elevada não é susceptível de ser estável durante as décadas de transição. Em vez disso, quanto mais elevada for a taxa, mais variável ela tende a ser. Isso tem sido demonstrado empiricamente por diferenças entre países nas últimas décadas (Jaffe e Kleiman 1975; Logue e Willett 1976). Também é altamente plausível em bases teóricas — tanto sobre a inflação real quanto, ainda mais claramente, as antecipações dos agentes econômicos com respeito à inflação. Os governos não têm a inflação alta como uma política deliberada anunciada, mas como consequência de outras políticas — em particular, políticas de pleno emprego e políticas de estado de bem-estar social que aumentam os gastos governamentais. Todos proclamam sua adesão ao objetivo de preços estáveis. Eles o fazem em resposta aos seus constituintes, que podem acolher muitos dos efeitos colaterais da inflação, mas ainda estão casados com o conceito de moeda estável. Uma explosão de inflação produz forte pressão para combatê-la. A política vai de uma direção para outra, encorajando ampla variação na inflação real e antecipada. E, claro, em tal ambiente, ninguém tem antecipações de valor único. Todos reconhecem que há grande incerteza sobre qual será a inflação real ao longo de qualquer intervalo futuro específico (Jaffe e Kleiman 1975; Meiselman 1976).

A tendência de que a inflação seja alta na média para ser altamente variável é reforçada pelo efeito da inflação na coesão política de um país no qual arranjos institucionais e contratos financeiros foram ajustados a um nível de preço "normal" de longo prazo. Alguns grupos ganham (por exemplo, proprietários de imóveis); outros perdem (por exemplo, proprietários de contas de poupança e títulos de renda fixa). O comportamento "prudente" torna-se, na verdade, imprudente, e o comportamento "imprudente" torna-se, na verdade, prudente. A sociedade é polarizada; um grupo é colocado contra outro. A agitação política aumenta. A capacidade de qualquer governo de governar é reduzida ao mesmo tempo em que a pressão por ações governamentais fortes cresce.

Uma maior variabilidade da inflação real ou antecipada pode elevar a taxa natural de desemprego de duas maneiras bastante diferentes.

Primeiro, o aumento da volatilidade encurta a duração ideal de compromissos não indexados e torna a indexação mais vantajosa (Gray 1976). Mas leva tempo para que a prática real se ajuste. Enquanto isso, arranjos anteriores introduzem rigidez que reduz a eficácia dos mercados. Um elemento adicional de incerteza é, por assim dizer, adicionado a cada arranjo de mercado. Além disso, a indexação é, mesmo na melhor das hipóteses, um substituto imperfeito para a estabilidade da taxa de inflação. Os índices de preços são imperfeitos; eles estão disponíveis apenas com um atraso e geralmente são aplicados aos termos do contrato apenas com um atraso adicional.

Estes desenvolvimentos claramente diminuem a eficiência econômica. É menos claro qual é o seu efeito no desemprego registrado. Inventários médios elevados de todos os tipos são uma forma de enfrentar a rigidez e a incerteza crescentes. Mas isso pode significar entesouramento de mão de obra pelas empresas e baixo desemprego ou uma força de trabalho maior entre empregos e, portanto, alto desemprego. Compromissos mais curtos podem significar um ajuste mais rápido do emprego às mudanças nas condições e, portanto, baixo desemprego, ou o atraso no ajuste da duração dos compromissos pode levar a um ajuste menos satisfatório e, portanto, alto desemprego. Claramente, muita pesquisa adicional é necessária nesta área para esclarecer a importância relativa dos vários efeitos. Tudo o que se pode dizer agora é que o ajuste lento dos compromissos e as imperfeições da indexação podem contribuir para o aumento registrado no desemprego.

Um segundo efeito relacionado ao aumento da volatilidade da inflação é tornar o sistema de preços um sistema menos eficiente para coordenar a atividade econômica. Uma função fundamental de um sistema de preços, como Hayek (1945) enfatizou tão brilhantemente, é transmitir de forma compacta, eficiente e a baixo custo a informação de que os agentes econômicos precisam para decidir o que produzir e como produzi-lo, ou como empregar recursos próprios. A informação relevante é sobre preços \textit{relativos} — o preço de um produto em relação a outro, das fontes de um fator de produção em relação a outro, de produtos em relação a serviços de fatores, de preços agora em relação a preços no futuro. Mas a informação na prática é transmitida na forma de preços \textit{absolutos} — preços em dólares ou libras ou coroas. Se o nível de preço estiver estável na média ou mudando a uma taxa constante, é relativamente fácil extrair o sinal sobre preços relativos dos preços absolutos observados. Quanto mais volátil a taxa de inflação geral, mais difícil se torna extrair o sinal sobre preços relativos dos preços absolutos: a transmissão de preços relativos é, por assim dizer, sendo congestionada pelo ruído vindo da transmissão de inflação (Lucas 1973, 1975; Harberger 1976). No extremo, o sistema de preços absolutos torna-se quase inútil, e os agentes econômicos recorrem a uma moeda alternativa ou à troca, com efeitos desastrosos na produtividade.

Novamente, o efeito sobre a eficiência econômica é claro, no desemprego menos. Mas, novamente, parece plausível que o nível médio de desemprego seria elevado pela maior quantidade de ruído nos sinais do mercado, pelo menos durante o período em que os arranjos institucionais não estão ainda adaptados à nova situação.

Estes efeitos do aumento da volatilidade da inflação ocorreriam mesmo se os preços fossem legalmente livres para se ajustar — se, nesse sentido, a inflação fosse aberta. Na prática, os efeitos distorcivos da incerteza, a rigidez de contratos voluntários de longo prazo e a contaminação dos sinais de preços seriam quase certamente reforçados por restrições legais na mudança de preços. No mundo moderno, os governos são eles mesmos produtores de serviços vendidos no mercado: de serviços postais a uma ampla gama de outros itens. Outros preços são regulados pelo governo e requerem aprovação governamental para mudança: de passagens aéreas a tarifas de táxi e cobranças de eletricidade. Nestes casos, os governos não podem evitar estar envolvidos no processo de fixação de preços. Além disso, as forças sociais e políticas desencadeadas por taxas de inflação voláteis levarão os governos a tentar reprimir a inflação em outras áreas: por controle explícito de preços e salários, ou por pressionar empresas ou sindicatos privados "voluntariamente" a exercer "contencão", ou por especular em câmbio para alterar a taxa de câmbio.

Os detalhes variarão de tempos em tempos e de país para país, mas o resultado geral é o mesmo: redução na capacidade do sistema de preços para guiar a atividade econômica; distorções nos preços relativos por causa da introdução de maior fricção, por assim dizer, em todos os mercados; e, muito provavelmente, uma maior taxa registrada de desemprego (Friedman 1976, cap. 12).

As forças que acabo de descrever podem tornar o sistema político e econômico dinamicamente instável e produzir hiperinflação e mudança política radical — como em muitos países derrotados após a Primeira Guerra Mundial, ou no Chile e na Argentina mais recentemente. No outro extremo, antes que tal catástrofe ocorra, políticas podem ser adotadas que alcançarão uma inflação relativamente baixa e estável e levarão ao desmantelamento de muitas das interferências com o sistema de preços. Isso restabeleceria as precondições para a hipótese da taxa natural direta e permitiria que essa hipótese fosse usada para prever o curso da transição.

Uma possibilidade intermediária é que o sistema alcance estabilidade a uma taxa média de inflação razoavelmente constante, embora alta. Nesse caso, o desemprego também deve se estabilizar em um nível razoavelmente constante, decididamente menor do que durante a transição. Como a discussão anterior enfatiza, a volatilidade crescente e a intervenção governamental crescente com o sistema de preços são os principais fatores que parecem prováveis de aumentar o desemprego, não a inflação alta ou uma taxa alta de intervenção por si só.

Maneiras de lidar com a volatilidade e a intervenção se desenvolverão: através da indexação e arranjos similares para lidar com a volatilidade da inflação; através do desenvolvimento de formas indiretas de alterar preços e salários para evitar controles governamentais.

Sob estas circunstâncias, a curva de Phillips de longo prazo seria novamente vertical e estaríamos de volta à hipótese da taxa natural, embora talvez para uma taxa de inflação natural diferente daquela para a qual ela foi sugerida pela primeira vez.

Como o fenômeno a ser explicado é a coexistência de inflação alta e desemprego alto, estressei o efeito das mudanças institucionais produzidas por uma transição de um sistema monetário no qual havia um nível de preço "normal" para um sistema monetário consistente com longos períodos de inflação alta e possivelmente altamente variável. Deve-se notar que uma vez que essas mudanças institucionais foram feitas, e os agentes econômicos ajustaram suas práticas e antecipações a elas, uma reversão para a estrutura monetária anterior ou mesmo a adoção de um novo quadro monetário de uma política bem-sucedida de baixa inflação exigiria, por sua vez, novos ajustes, e estes poderiam ter muitos dos mesmos efeitos transicionais adversos no nível de emprego. Haveria o que parece ser uma curva de Phillips de médio prazo negativamente inclinada em vez da positivamente inclinada que tentei racionalizar.

\section{Conclusão}

Uma consequência da revolução Keynesiana dos anos 1930 foi a aceitação de um nível de salário absoluto rígido, e um nível de preço absoluto quase rígido, como ponto de partida para analisar a mudança econômica de curto prazo. Passou-se a considerar que estes eram dados essencialmente institucionais e, portanto, eram considerados por agentes econômicos, de modo que as mudanças na demanda nominal agregada seriam refletidas quase inteiramente na produção e dificilmente em todos os preços. A confusão milenar entre preços absolutos e preços relativos ganhou uma nova sobrevida.

Neste ambiente intelectual, era compreensível que os economistas analisassem a relação entre desemprego e demanda nominal em vez de salários reais e considerassem implicitamente as mudanças na demanda nominal como iguais às mudanças nos salários reais antecipados. Além disso, a evidência empírica que inicialmente sugeriu uma relação estável entre o nível de desemprego e a taxa de mudança de salários nominais foi extraída de um período em que, apesar de flutuações acentuadas de curto prazo nos preços, havia um nível de preço estável a longo prazo e quando a expectativa de estabilidade contínua era amplamente compartilhada. Daí que esses dados exibissem nenhum sinal de aviso sobre o caráter especial das premissas.

A hipótese de que existe uma relação estável entre o nível de desemprego e a taxa de inflação foi adotada pela profissão econômica com alacridade. Preencheu uma lacuna na estrutura teórica de Keynes. Parecia ser a "equação única" que o próprio Keynes dissera que "estamos... em falta" (1936, p. 276). Além disso, parecia fornecer uma ferramenta confiável para a política econômica, permitindo que o economista informasse o formulador de políticas sobre as alternativas disponíveis para ele.

Em qualquer ciência, contanto que a experiência parecesse ser consistente com a hipótese reinante, ela continuou a ser aceita, embora, como sempre, alguns poucos dissidentes questionassem sua validade. Mas à medida que os anos 50 se tornaram os anos 60, e os anos 60 os anos 70, tornou-se cada vez mais difícil aceitar a hipótese em sua forma simples. Parecia requerer doses cada vez maiores de inflação para manter baixo o nível de desemprego. A estagflação ergueu sua cabeça feia.

Muitas tentativas foram feitas para remendar a hipótese, permitindo fatores especiais, como a força dos sindicatos. Mas a experiência teimosamente recusou-se a conformar-se às versões remendadas.

Uma revisão mais radical era necessária. Tomou a forma de enfatizar a importância das surpresas — de diferenças entre magnitudes reais e antecipadas. Restaurou a primazia da distinção entre magnitudes reais e nominais. Existe uma taxa natural de desemprego a qualquer tempo determinada por fatores reais. Esta taxa natural será atingida quando as expectativas forem, na média, realizadas. A mesma situação real é consistente com qualquer nível absoluto de preços ou de mudança de preços, desde que seja feita provisão para o efeito da mudança de preço no custo real de manutenção de saldos monetários. Nesse aspecto, a moeda é neutra. Por outro lado, mudanças não antecipadas na demanda nominal agregada e na inflação causarão erros sistemáticos de percepção por parte dos empregadores e empregados que inicialmente levarão o desemprego a desviar-se na direção oposta de sua taxa natural. Nesse aspecto, a moeda não é neutra. No entanto, tais desvios são transitórios, embora possa levar um longo tempo cronológico antes que sejam revertidos e finalmente eliminados à medida que as antecipações se ajustam.

A hipótese da taxa natural contém a curva de Phillips original como um caso especial e racionaliza uma gama muito mais ampla de experiências, em particular o fenômeno da estagflação. Ela agora foi amplamente, embora não universalmente, aceita.

No entanto, a hipótese da taxa natural em sua forma atual não provou ser rica o suficiente para explicar um desenvolvimento mais recente — um movimento de estagflação para "slumpflation" [recessão com inflação]. Em anos recentes, a inflação mais alta foi frequentemente acompanhada por um desemprego mais elevado, não mais baixo, como a curva de Phillips simples sugeriria, nem o mesmo desemprego, como a hipótese da taxa natural sugeriria.

Esta associação recente de inflação mais alta com desemprego mais alto pode refletir o impacto comum de eventos como a crise do petróleo, ou forças independentes que conferiram uma tendência ascendente comum à inflação e ao desemprego.

No entanto, um fator importante em alguns países e um fator contribuinte em outros pode ser que eles estão em um período de transição — este tempo a ser medido por quinquênios ou décadas e não por anos. O público não adaptou suas atitudes ou suas instituições a um novo ambiente monetário. A inflação tende não apenas a ser mais alta, mas também cada vez mais volátil e a ser acompanhada por uma intervenção governamental crescente no estabelecimento de preços. A volatilidade crescente da inflação e a partida crescente dos preços relativos dos valores que as forças de mercado sozinhas estabeleceriam combinam-se para tornar o sistema econômico menos eficiente, para introduzir fricções em todos os mercados e, muito provavelmente, para aumentar a taxa registrada de desemprego.

Nesta análise, a situação atual não pode durar. Ou ela degenerará em hiperinflação e mudança radical, ou as instituições se ajustarão a uma situação de inflação crônica, ou os governos adotarão políticas que produzirão uma taxa baixa de inflação e menos intervenção governamental na fixação de preços.

Contei uma história perfeitamente padrão de como as teorias científicas são revisadas. No entanto, é uma história que tem uma importância de longo alcance. O governo e a política sobre inflação e desemprego têm estado no centro da controvérsia política. A guerra ideológica tem assolado estas questões. No entanto, a mudança drástica que ocorreu na teoria econômica não foi o resultado de uma guerra ideológica. Não resultou de crenças ou objetivos políticos divergentes. Ela respondeu quase inteiramente à força dos eventos: a experiência bruta provou ser muito mais potente do que a mais forte das preferências políticas ou ideológicas.

A importância para a humanidade de uma compreensão correta da ciência econômica positiva é vividamente trazida à tona por uma declaração feita há quase duzentos anos por Pierre S. du Pont, um deputado de Nemours para a Assembleia Nacional Francesa, falando, apropriadamente, sobre uma proposta para emitir \textit{assignats} adicionais — a moeda fiduciária da Revolução Francesa: "Senhores, é um costume desagradável para o qual se é levado muito facilmente pela dureza das discussões, assumir intenções más. É necessário ser gracioso quanto às intenções; deve-se acreditar que elas são boas, e aparentemente são; mas não temos que ser graciosos de forma alguma para a lógica inconsistente ou para o raciocínio absurdo. Maus lógicos cometeram crimes mais involuntários do que homens maus fizeram intencionalmente" (25 de setembro de 1790).

\section*{Referências}

Friedman, Milton. "The Methodology of Positive Economics." In \textit{Essays in Positive Economics}. Chicago: Univ. Chicago Press, 1953.

———. "An Inflationary Recession." \textit{Newsweek} (17 de outubro de 1966). (\textit{a})

———. "What Price Guideposts?" In \textit{Guidelines: Informal Contracts and the Market Place}, editado por G. P. Shultz e R. Z. Aliber. Chicago: Univ. Chicago Press, 1966. (\textit{b})

———. \textit{Inflation: Causes and Consequences}. Bombay: Asia Publishing House, 1963. Republicado em \textit{Dollars and Deficits}. Englewood Cliffs, N.J.: Prentice-Hall, 1968. (\textit{a})

———. "The Role of Monetary Policy." \textit{A.E.R.} 58 (Março de 1968): 1–17. (\textit{b})

———. \textit{Price Theory}. Chicago: Aldine, 1976.

Gordon, Robert J. "Alternative Responses of Policy to External Supply Shocks." \textit{Brookings Papers Econ. Activity}, no. 1 (1975). (\textit{a})

———. "The Demand and Supply of Inflation." \textit{J. Law and Econ.} 18 (Dezembro de 1975): 807–36. (\textit{b})

———. "Recent Developments in the Theory of Inflation and Unemployment." \textit{J. Monetary Econ.} 2 (1976): 185–219.

Gray, Jo Anna. "Essays on Wage Indexation." Ph.D. dissertation, Univ. Chicago, 1976.

Harberger, Arnold G. "The Inflation Problem in Latin America." Um relatório preparado para o encontro em Buenos Aires (Março de 1966) do Comitê Interamericano da Aliança para o Progresso, publicado em espanhol como "El problema de la inflación en America Latina," in \textit{Centro de Estudios Monetarios Latinoamericanos, Boletin mensual} (Junho 1966), pp. 253–69. Republicado em \textit{Economic Development Institute, Trabajos sobre desarrollo económico}. Washington: IBRD, 1967.

———. "Inflation." In \textit{The Great Ideas Today}, 1976, editado por Robert M. Hutchins e Mortimer J. Adler. Chicago: Encyclopaedia Britannica, 1976.

Hayek, F. A. "The Use of Knowledge in Society." \textit{A.E.R.} 35 (Setembro de 1945): 519–30.

Jaffe, Dwight, e Kleiman, Ephraim. "The Welfare Implications of Uneven Inflation." Seminar paper no. 50, Inst. Internat. Econ. Studies, Univ. Stockholm, Novembro 1975.

Keynes, J. M. \textit{General Theory of Employment, Interest, and Money}. Londres: Macmillan, 1936.

Klein, Benjamin. "Our New Monetary Standard: The Measurement and Effects of Price Uncertainty, 1880–1973." \textit{Econ. Inquiry} 13 (Dezembro 1975): 461–84.

Logue, Dennis E., e Willett, Thomas D. "A Note on the Relation between the Rate and Variability of Inflation." \textit{Economica} 43 (Maio 1976): 151–58.

Lucas, Robert E. "Some International Evidence on Output-Inflation Trade-offs." \textit{A.E.R.} 63 (Junho 1973): 326–34.

———. "An Equilibrium Model of the Business Cycle." \textit{J.P.E.} 83, no. 6 (Dezembro 1975): 1113–44.

Meiselman, David. "Capital Formation, Monetary and Financial Adjustments." \textit{Proceedings, 27th National Conference of Tax Foundation}. Washington: Tax Foundation, 1976.

Muth, John. "Rational Expectations and the Theory of Price Movements." \textit{Econometrica} 29 (Julho 1961): 315–33.

Phelps, E. S. "Phillips Curve, Expectations of Inflation and Optimal Unemployment over Time." \textit{Economica}, n.s. 34 (Agosto 1967): 254–81.

———. "Money Wage Dynamics and Labor Market Equilibrium." In \textit{Microeconomic Foundations of Employment and Inflation Theory}, editado por E. S. Phelps. Nova York: Norton, 1970.

———. "The Relationship between Unemployment and the Rate of Change of Money Wage Rates in the United Kingdom, 1861–1957." \textit{Economica} (Novembro 1958), pp. 283–99.

Sjaastad, Larry A. "Monetary Policy and Suppressed Inflation in Latin America." In \textit{National Monetary Policies and the International Financial System}, editado por Robert Z. Aliber. Chicago: Univ. Chicago Press, 1974.

\end{document}